% Blueprint for the formalization of
%   L. Becchetti, A. Clementi, E. Natale, F. Pasquale, R. Silvestri, L. Trevisan,
%   "Simple Dynamics for Plurality Consensus", SPAA 2014 (arXiv:1310.2858).
% Numbering follows the SPAA 2014 proceedings version. Every statement carries
% \lean{...} / \leanok / \uses{...} tags matched to the Lean declarations in
% Plurality/*.lean (namespace `Plurality`). Statements without \leanok are
% not formalized yet. Where the Lean proof departs from the paper, this file
% follows the Lean proof; the departures are listed in
% FORMALIZATION_DIFFERENCES.md.

\begin{abstract}
\noindent
$n$ anonymous nodes each support one of $k$ colors. In every round each node
samples three nodes uniformly at random (with repetition, possibly itself)
and adopts the majority color among them, or the first one if all three
differ. If some color $m$ leads every other color by a sufficiently large
\emph{bias} $s$, the process reaches consensus on $m$. The paper proves that
this takes $O(\lambda\log n)$ rounds w.h.p.\ whenever $c_m \ge n/\lambda$
and $s \ge c\sqrt{\lambda n\log n}$ (Theorem~\ref{thm:3_8}), that
$\Omega(k\log n)$ rounds are needed from balanced configurations, that
3-majority is essentially the only 3-input rule that solves plurality
consensus, and that sampling $h > 3$ nodes speeds the process up by at most
a factor $O(h^2)$.

The upper bound (Section~\ref{sec:upper}), including the headline bound
$O(\min\{k,(n/\log n)^{1/3}\}\log n)$ of Corollary~\ref{cor:3_10}, is fully
formalized, with no
\texttt{sorry}, on the finite probability layer of the shared
\texttt{dynamics} package: all randomness is uniform over finite types,
multi-round probabilities are iterates of a finite Markov kernel, and the
concentration inequalities (Chernoff, Hoeffding, Bernstein) are proved from
scratch. The lower bounds of Section~\ref{sec:lower} are formalized as well,
with two exceptions where the paper's argument does not go through: the
$\Omega(k\log n)$ bound is proved for $k \le n^{1/4-\delta}$ rather than up to
$(n/\log n)^{1/4}$, and part~(a) of Theorem~4.8 (Theorem~\ref{thm:4_8_a}) leaves open rules
with $\Delta_r, \Delta_b \le 1$. Observation~3.9 (an adversary) is not
formalized. Beyond the paper, Section~\ref{sec:any_start} proves that with two opinions
3-majority reaches consensus from any configuration, including a perfectly balanced one, within
$O(\log n)$ rounds w.h.p.
\end{abstract}

\section{The model}\label{sec:model}

\begin{definition}[Colorings and color counts]\label{def:config}
  \lean{Plurality.Config, Plurality.count, Plurality.sum_count}
  \leanok
  A \emph{coloring} is a map $x \colon [n] \to [k]$. The number of nodes of
  color $j$ is $c_j = |\{v : x(v) = j\}|$, and $\sum_j c_j = n$. The vector
  $c = (c_1,\dots,c_k)$ is the paper's $k$-color distribution ($k$-cd).
\end{definition}

\begin{definition}[3-input rules and the 3-majority rule]\label{def:rule}
  \lean{Plurality.maj3, Plurality.maj3_mem}
  \leanok
  A \emph{3-input rule} is a map $f \colon [k]^3 \to [k]$ with
  $f(a,b,c) \in \{a,b,c\}$. The \emph{3-majority rule} is
  $\mathrm{maj}_3(a,b,c) = b$ if $b = c$ and $a$ otherwise; it returns the
  majority color, and the first color if all three differ.
\end{definition}

\begin{definition}[One round and the process]\label{def:step}
  \lean{Plurality.stepWith, Plurality.step, Plurality.runWith, Plurality.run}
  \leanok
  \uses{def:config, def:rule}
  A round is a map $r \colon [n] \to [n]^3$ (the three samples of every node),
  drawn uniformly. The next coloring under rule $f$ is
  $v \mapsto f(x(r_v^1), x(r_v^2), x(r_v^3))$. The process after $T$ rounds
  folds this step over $T$ independent uniform rounds.
\end{definition}

\begin{definition}[Consensus]\label{def:mono}
  \lean{Plurality.Mono, Plurality.mono_iff_count, Plurality.stepWith_mono, Plurality.run_mono}
  \leanok
  \uses{def:step}
  $x$ is in \emph{consensus on $m$} if every node has color $m$, i.e.\
  $c_m = n$. For every 3-input rule, consensus is absorbing.
\end{definition}

\begin{definition}[The Markov kernel]\label{def:kernel}
  \lean{Plurality.kernel, Plurality.kernel_event, Dynamics.Kernel.ofStep, Dynamics.Kernel.iterate_ofStep}
  \leanok
  \uses{def:step}
  The process with rule $f$ is the finite Markov kernel on colorings whose
  row at $x$ is the law of one round of $f$ from $x$. Its $T$-step event
  probabilities are the expectations over $T$ independent uniform rounds.
\end{definition}

\section{The next expected coloring}\label{sec:expected}

\begin{lemma}[Adoption probability]\label{lem:adopt}
  \lean{Plurality.count_stepWith_eq_sum, Plurality.avg_adopt, Plurality.sum_triple_maj3, Plurality.avg_adopt_maj3}
  \leanok
  \uses{def:step}
  The new count of color $j$ is a sum of $n$ independent indicators, one per
  node. Under 3-majority, a node adopts $j$ with probability
  $x_j^2 + x_j(1 - \sum_h x_h^2)$, where $x_h = c_h/n$.
\end{lemma}

\begin{proof}
  The three samples are independent with law $(x_h)_h$. Summing the
  indicator of $\mathrm{maj}_3(a,b,c) = j$ first over $c$, then $b$, then $a$
  gives the polynomial.
\end{proof}

\begin{lemma}[Lemma 2.1, next expected coloring]\label{lem:2_1}
  \lean{Plurality.expected_count}
  \leanok
  \uses{lem:adopt}
  For every coloring and every color $j$,
  \[
    \mu_j(c) = \E[C_{j,t+1} \mid C_t = c]
      = c_j\Bigl(1 + \frac{1}{n^2}\Bigl(n c_j - \sum_h c_h^2\Bigr)\Bigr).
  \]
\end{lemma}

\begin{proof}
  Multiply the adoption probability of Lemma~\ref{lem:adopt} by $n$.
\end{proof}

\section{The upper bound}\label{sec:upper}

This section formalizes Theorem~3.8 of Becchetti, Clementi, Natale, Pasquale, Silvestri and
Trevisan, \emph{Simple Dynamics for Plurality Consensus} (SPAA 2014), and its Corollaries 3.10
to~3.12. The SPAA proof of Theorem~3.8 needs a minor correction (small repairs: the saturation
phases as defined are not shown to be preserved, and a union bound over the competing colors is
missing), and the repaired proof is what is formalized
(\href{https://github.com/formal-dynamics/leanamics/blob/main/plurality/FORMALIZATION_DIFFERENCES.md#1-theorem-38-how-the-lean-proof-differs-from-spaa}{details}).

\subsection{Plurality, bias, \texorpdfstring{$\alpha$}{alpha} and \texorpdfstring{$\gamma$}{gamma}}

\begin{definition}[Section 3 quantities]\label{def:quantities}
  \lean{Plurality.maxc, Plurality.argmaxSet, Plurality.secondc, Plurality.bias, Plurality.alpha, Plurality.gamma, Plurality.mu}
  \leanok
  For a $k$-cd $c$ with $\sum_h c_h = n$: $m(c) = \max_h c_h$,
  $M(c) = \{j : c_j = m(c)\}$, the \emph{bias}
  $s(c) = m(c) - \max_{h \notin M(c)} c_h$ if $|M(c)| = 1$ and $0$ otherwise,
  $\alpha(c) = (n - m(c))\,s(c)/n^2$ and
  $\gamma(c) = (n\,m(c) - \sum_h c_h^2)/n^2 - \alpha(c)$.
\end{definition}

\begin{lemma}[Every other color trails by the bias]\label{lem:bias_gap}
  \lean{Plurality.add_bias_le, Plurality.exists_add_bias_eq, Plurality.gamma_mul_sq}
  \leanok
  \uses{def:quantities}
  If $c_m = m(c)$ then $c_h + s(c) \le m(c)$ for every $h \ne m$, with
  equality for some $h$ when $k \ge 2$. Moreover
  $\gamma(c)\,n^2 = \sum_{h \ne m} c_h\,(m(c) - s(c) - c_h)$.
\end{lemma}

\begin{lemma}[Lemma 3.1]\label{lem:3_1}
  \lean{Plurality.lemma_3_1_a, Plurality.lemma_3_1_b, Plurality.lemma_3_1_c}
  \leanok
  \uses{lem:bias_gap}
  For every $k$-cd $c$:
  (a) $0 \le s(c) \le m(c) - (n - m(c))/(k-1)$;
  (b) $0 \le \alpha(c) \le \min\{s(c)/n,\ 1/4\}$;
  (c) $0 \le \gamma(c) \le 1/8$.
\end{lemma}

\begin{proof}
  (a) The $k-1$ other colors each have at most $m(c) - s(c)$ nodes.
  (b) $(n - m)m \le n^2/4$ and $s \le m$.
  (c) Use the identity of Lemma~\ref{lem:bias_gap}: every summand is
  nonnegative, the color attaining the bias contributes $0$, and the rest is
  at most $d(n - m - d) \le n^2/8$ with $d = m(c) - s(c)$.
\end{proof}

\begin{lemma}[Lemma 3.2]\label{lem:3_2}
  \lean{Plurality.lemma_3_2_a, Plurality.lemma_3_2_b, Plurality.lemma_3_2_c}
  \leanok
  \uses{def:quantities, lem:bias_gap}
  Let $m \in M(c)$ and let $\ell \ne m$ have the most nodes among the other
  colors. Then (a) $\mu_m(c) = c_m(1 + \gamma + \alpha)$;
  (b) $\mu_j(c) = c_j(1 + \gamma + \alpha - (m(c) - c_j)/n)$ for every $j$;
  (c) $\mu_\ell(c) = c_\ell(1 + \gamma + \alpha - s(c)/n)$.
\end{lemma}

\subsection{Concentration}

\begin{lemma}[Concentration for sums of independent coordinates]\label{lem:tails}
  \lean{Dynamics.avg_bernstein, Dynamics.avg_hoeffding, Dynamics.avg_chernoff_lower_mul,
    Dynamics.avg_markov_one, Dynamics.avg_chernoff_upper_log}
  \leanok
  For $X = \sum_i Y_i(\omega_i)$ with independent uniform coordinates:
  Bernstein's inequality
  $\Pr(X \ge \E X + \lambda) \le \exp(-\lambda^2/(2\sigma^2(1 + b\lambda/(3\sigma^2))))$
  when $Y_i - \E Y_i \le b$ and $\sum_i \mathrm{Var}\,Y_i \le \sigma^2$;
  Hoeffding's inequality for $\{0,1\}$ coordinates; the multiplicative
  Chernoff lower tail $\Pr(X \le (1-\delta)\mu) \le \exp(-\delta^2\mu/2)$;
  Markov's inequality $\Pr(X \ge 1) \le \E X$ for nonnegative coordinates;
  and the closed-form Chernoff upper tail
  $\Pr(X \ge k) \le \exp(k - \mu - k\log(k/\mu))$ for $\E X \le \mu \le k$.
  All are proved once in the shared \texttt{dynamics} library
  (\texttt{Dynamics/Concentration.lean}, \texttt{Dynamics/Chernoff.lean},
  \texttt{Dynamics/Tail.lean}).
\end{lemma}

\subsection{Growth of the bias}

\begin{lemma}[Lemma 3.3, increasing rate of the bias]\label{lem:3_3}
  \lean{Plurality.lemma_3_3}
  \leanok
  \uses{lem:2_1, lem:3_1, lem:3_2, lem:tails}
  If $M(c) = \{m\}$ then for every $j \ne m$,
  \[
    \Pr\Bigl(C_{m,t+1} - C_{j,t+1} \le s(c)\Bigl(1 + \gamma(c)
      + \frac{c_m\alpha(c)}{2s(c)}\Bigr) \Bigm| C_t = c\Bigr)
      \le \exp\Bigl(-\frac{c_m\alpha(c)^2}{25}\Bigr).
  \]
\end{lemma}

\begin{proof}
  By Lemmas~\ref{lem:3_1} and~\ref{lem:3_2},
  $\E[C_{m,t+1} - C_{j,t+1}] \ge s(1+\gamma) + c_m\alpha$. The difference is
  a sum of per-node variables in $\{-1,0,1\}$ with variance at most $3c_m$;
  apply Bernstein's inequality with $b = 2$ and $\lambda = c_m\alpha/2$.
\end{proof}

\begin{lemma}[Lemma 3.4]\label{lem:3_4}
  \lean{Plurality.lemma_3_4}
  \leanok
  \uses{lem:3_1, lem:3_2, lem:tails}
  If $M(c) = \{m\}$ then
  $\Pr(C_{m,t+1} \le c_m(1 + \gamma + \alpha/2) \mid C_t = c)
  \le \exp(-c_m\alpha^2/11)$.
\end{lemma}

\begin{proof}
  The multiplicative Chernoff lower tail with
  $\delta = \alpha/(2(1 + \gamma + \alpha))$.
\end{proof}

\begin{lemma}[Lemma 3.5, large plurality and large bias]\label{lem:3_5}
  \lean{Plurality.lemma_3_5}
  \leanok
  \uses{lem:3_3, lem:3_4}
  If $M(c) = \{m\}$, $0 < \lambda \le 2/3$, $\lambda n \le c_m \le (2/3)n$
  and $s(c) \ge 22\sqrt{(1/\lambda)\,n\log n}$, then for every $j \ne m$
  \[
    \Pr\bigl(C_{m,t+1} - C_{j,t+1} \le s(c)(1 + \lambda/6)\bigr) \le 1/n^2
    \quad\text{and}\quad
    \Pr\bigl(C_{m,t+1} \le c_m\bigr) \le 1/n^2.
  \]
\end{lemma}

\begin{proof}
  Under the hypotheses $c_m\alpha/(2s) \ge \lambda/6$ and
  $c_m\alpha^2/25 \ge 2\log n$; conclude with Lemmas~\ref{lem:3_3}
  and~\ref{lem:3_4}.
\end{proof}

\subsection{Saturation}

\begin{lemma}[Lemma 3.6]\label{lem:3_6}
  \lean{Plurality.lemma_3_6}
  \leanok
  \uses{lem:2_1, lem:3_1}
  For every $m \in M(c)$, writing $\bar C_{m,t+1} = n - C_{m,t+1}$,
  \[
    (n - c_m)\Bigl(1 - \frac{c_m^2}{n^2}\Bigr)
      \le \E[\bar C_{m,t+1} \mid C_t = c]
      \le (n - c_m)\Bigl(1 - \frac{s(c)\,c_m}{n^2}\Bigr).
  \]
\end{lemma}

\begin{lemma}[Lemma 3.7, very large plurality]\label{lem:3_7}
  \lean{Plurality.lemma_3_7_i, Plurality.lemma_3_7_ii, Plurality.bias_ge_of_singleton}
  \leanok
  \uses{lem:3_6, lem:tails}
  Let $\log n \ge 40$, $M(c) = \{m\}$.
  (i) If $s(c) \ge n/3$ and $n - c_m \ge n^{1/4}\log n$ then
  $\Pr(\bar C_{m,t+1} \ge \frac{17}{18}(n - c_m)) \le 1/n^2$.
  (ii) If $n - c_m < n^{1/4}\log n$ then $\Pr(\bar C_{m,t+1} > 0) \le n^{-1/5}$
  and $\Pr(\bar C_{m,t+1} \ge n^{1/4}\log n) \le 1/n^2$.
\end{lemma}

\begin{proof}
  The paper omits this proof. (i) By Lemma~\ref{lem:3_6} and
  $s\,c_m \ge n^2/9$, $\E\bar C_{m,t+1} \le \frac89(n - c_m)$; apply the
  closed-form Chernoff upper tail at ratio $17/16$. (ii) Here
  $s(c) \ge n - 2(n - c_m)$, so $\E\bar C_{m,t+1} \le 3(n-c_m)^2/n
  \le 3\log^2 n\, e^{-(\log n)/2}$; apply Markov's inequality for the first
  bound and the Chernoff upper tail for the second.
\end{proof}

\subsection{Nested phases and the main theorem}

\begin{lemma}[Lemma A.4, nested phases]\label{lem:A_4}
  \lean{Dynamics.Kernel.nested_phases}
  \leanok
  Let $A_1 \supseteq \cdots \supseteq A_T$ be sets of states of a finite
  Markov chain such that from every state of $A_i$ the chain stays in $A_i$
  with probability $\ge 1-\varepsilon$ and, for $i < T$, moves into
  $A_{i+1}$ with probability $\ge 1 - \nu$. Then from any state of $A_1$,
  after $\ell T$ steps the chain is in $A_T$ with probability at least
  $1 - T(\ell\varepsilon + \nu^\ell)$.
\end{lemma}

\begin{proof}
  The paper omits this proof. Phase by phase: within $\ell$ steps the chain
  leaves $A_i$ with probability at most $\ell\varepsilon$, and fails to jump
  into $A_{i+1}$ in each of the $\ell$ steps with probability at most
  $\nu^\ell$. The hypothesis $\varepsilon \le \nu$ of the paper is not needed.
\end{proof}

\begin{definition}[The phases]\label{def:phases}
  \lean{Plurality.growthSet, Plurality.satSet, Plurality.monoSet}
  \leanok
  \uses{def:quantities, def:mono}
  Fix $m$, $\lambda$, $\Lambda$ and let $q = 1 + 1/(6\lambda)$,
  $B = n^{1/4}\log n$.
  \emph{Growth phase $i$}: $c_m > \frac23 n$, or $M(c) = \{m\}$,
  $c_m \ge n/\lambda$ and $s(c) \ge \Lambda q^i$.
  \emph{Saturation phase $j$}:
  $n - c_m < \max\{\frac n3 (\frac{17}{18})^j,\ B\}$.
  \emph{Consensus}: $c_m = n$.
\end{definition}

\begin{remark}
  The paper's saturation phases also require $s(c) \ge n/3$, which
  Lemma~\ref{lem:3_7} does not show to be preserved. Here the saturation
  phases only bound $n - c_m < n/3$, which already implies $M(c) = \{m\}$
  and $s(c) \ge n/3$ (Lemma~\ref{lem:dissent}).
\end{remark}

\begin{lemma}[A clear majority]\label{lem:dissent}
  \lean{Plurality.of_dissent_lt, Plurality.argmaxSet_eq_of_lt, Plurality.le_bias_of_gap, Plurality.count_step_eq_zero}
  \leanok
  \uses{def:quantities}
  If $n - c_m < n/3$ then $M(c) = \{m\}$ and $s(c) \ge n/3$. If $m$ leads
  every other color by at least $g > 0$ and $c_m \ge g$, then $M(c) = \{m\}$
  and $s(c) \ge g$. A color without nodes stays without nodes.
\end{lemma}

\begin{lemma}[One round of each phase]\label{lem:phase_steps}
  \lean{Plurality.growth_step, Plurality.sat_step, Plurality.mono_step, Plurality.mono_stay, Plurality.growthSet_succ_sub, Plurality.satSet_succ_sub, Plurality.growthSet_sub_of_lt, Plurality.satSet_zero_sub}
  \leanok
  \uses{def:phases, def:kernel, lem:3_5, lem:3_7, lem:dissent}
  Let $\log n \ge 40$, $k \ge 2$, $\lambda \ge 3$ and
  $\Lambda \ge 22\sqrt{\lambda n \log n}$.
  From growth phase $i$ the chain enters growth phase $i+1$ with probability
  $\ge 1 - 1/n$. From saturation phase $j$ it enters saturation phase $j+1$
  with probability $\ge 1 - 1/n^2$. From fewer than $B$ dissenting nodes it
  reaches consensus with probability $\ge 1 - n^{-1/5}$, and consensus is
  kept with probability $1$. The phase sets are nested, and growth phase $i$
  lies in saturation phase $0$ once $\Lambda q^i > n$.
\end{lemma}

\begin{proof}
  \emph{Growth, case $c_m > \frac23 n$}: saturation phase $0$ moves into
  saturation phase $1$, which lies in $\{c_m > \frac23 n\}$.
  \emph{Growth, otherwise}: apply Lemma~\ref{lem:3_5} with $1/\lambda$ to each
  other color $j$ that currently has a node, and to the growth of $c_m$. There
  are fewer than $n - c_m$ such colors, so a union bound gives failure
  probability at most $n/n^2$. Colors without nodes stay empty, and the color
  attaining the bias has a node (otherwise $c_m = n$), so after a successful
  round $m$ leads every other color by $s(c)\,q$ and
  Lemma~\ref{lem:dissent} applies.
  \emph{Saturation}: Lemma~\ref{lem:3_7}(i) above $B$ dissenters and
  Lemma~\ref{lem:3_7}(ii) below.
\end{proof}

\begin{lemma}[Number of phases]\label{lem:phases}
  \lean{Plurality.growthPhases, Plurality.satPhases, Plurality.phases, Plurality.phases_le}
  \leanok
  \uses{def:phases}
  With $T_1 = \lceil \log n/\log q\rceil$ growth phases,
  $T_2 = \lceil \log n/\log(18/17)\rceil + 1$ saturation phases and one
  consensus phase, $T = T_1 + T_2 + 1 \le 13\lambda\log n$ for $\lambda \ge 3$
  and $\log n \ge 40$.
\end{lemma}

\begin{proof}
  $\log(1+x) \ge x/(1+x)$ gives $\log q \ge 1/(6\lambda+1)$ and
  $\log(18/17) \ge 1/18$.
\end{proof}

\begin{theorem}[Theorem 3.8, the general upper bound]\label{thm:3_8}
  \lean{Plurality.theorem_3_8, Plurality.theorem_3_8_bigO}
  \leanok
  \uses{lem:A_4, lem:phase_steps, lem:phases}
  Let $\log n \ge 40$, $k \ge 2$ and $\lambda \ge 3$. If $M(c) = \{m\}$,
  $c_m \ge n/\lambda$ and $s(c) \ge 22\sqrt{\lambda n\log n}$, then after
  $10T \le 130\lambda\log n$ rounds all nodes support $m$ with probability
  at least $1 - 11T/n \ge 1 - 143\lambda\log n/n$.
\end{theorem}

\begin{proof}
  Apply Lemma~\ref{lem:A_4} to the kernel of 3-majority with the $T_1$
  growth phases, the $T_2$ saturation phases and consensus,
  $\varepsilon = 1/n$, $\nu = n^{-1/5}$ and $\ell = 10$, so that
  $\nu^\ell = n^{-2}$. The initial coloring is in growth phase $0$; growth
  phase $T_1$ lies in saturation phase $0$ because $\Lambda q^{T_1} \ge 22n$;
  saturation phase $T_2 - 1$ lies below $B$ dissenters because
  $\frac n3(\frac{17}{18})^{T_2-1} \le \frac13$.
\end{proof}

\begin{remark}
  The paper states the failure probability as $11T/n^2$; the extra factor
  $n$ comes from the union bound over competing colors in
  Lemma~\ref{lem:phase_steps}, which the paper omits. The paper's
  hypothesis $\lambda < \sqrt n$ is not used; it only makes the hypothesis
  on $s(c)$ satisfiable.
\end{remark}

\begin{corollary}[Corollary 3.10]\label{cor:3_10}
  \lean{Plurality.corollary_3_10, Plurality.lamK, Plurality.le_mul_maxc}
  \leanok
  \uses{thm:3_8}
  If $M(c) = \{m\}$ and
  $s(c) \ge 22\sqrt{\min\{2k, (n/\log n)^{1/3}\}\,n\log n}$, then 3-majority
  reaches consensus on $m$ in $O(\min\{2k, (n/\log n)^{1/3}\}\log n)$
  rounds w.h.p. Precisely, for $\log n \ge 40$ and $k \ge 2$: at most
  $130\lambda\log n$ rounds with probability at least
  $1 - 143\lambda\log n/n$.
\end{corollary}

\begin{proof}
  \leanok
  Take $\lambda = \min\{2k, (n/\log n)^{1/3}\}$. If $\lambda = 2k$ then
  $c_m \ge n/k$ by pigeonhole; otherwise $\lambda^3 = n/\log n$, so
  $22\sqrt{\lambda n\log n} = 22n/\lambda$ and $c_m \ge s(c) \ge n/\lambda$.
  Moreover $\lambda \ge 3$ because $k \ge 2$ and $n = e^{\log n} \ge 27\log n$.
\end{proof}

\begin{corollary}[Corollaries 3.11 and 3.12]\label{cor:3_11}
  \lean{Plurality.corollary_3_11, Plurality.corollary_3_12}
  \leanok
  \uses{thm:3_8}
  If $c_m \ge n/\log^\ell n$ and $s(c) \ge 22\sqrt{n\log^{\ell+1}n}$,
  consensus takes $O(\log^{\ell+1} n)$ rounds w.h.p. If $c_m \ge n/\beta$
  and $s(c) \ge 22\sqrt{\beta n\log n}$ for a constant $\beta \ge 3$, it
  takes $O(\log n)$ rounds w.h.p.
\end{corollary}

\begin{proof}
  \leanok
  Take $\lambda = \log^\ell n$ (with $\ell \ge 1$, so that $\lambda \ge 3$)
  and $\lambda = \beta$ in Theorem~\ref{thm:3_8}.
\end{proof}

\subsection{Two colors}

\begin{theorem}[Binary case]\label{thm:binary}
  \lean{Plurality.colorSet_step, Plurality.colorSet_run, Plurality.binary_consensus_whp}
  \leanok
  \uses{def:step}
  With $k = 2$ the set of nodes of color $1$ evolves exactly as the binary
  3-majority process of the \texttt{three\_majority} package. Hence, if
  $c_1 \ge \frac35 n$ and $\log n \ge 30$, all nodes support color $1$ after
  $O(\log n)$ rounds with probability at least $1 - 500/n$.
\end{theorem}

\begin{corollary}[Two opinions from a vanishing bias]\label{cor:binary_vanishing}
  \lean{Plurality.ofSet, Plurality.colorSet_ofSet, Plurality.count_ofSet_one, Plurality.count_ofSet_zero, Plurality.majority3_vanishing_bias}
  \leanok
  \uses{thm:3_8, thm:binary}
  In the binary 3-majority process, if the initial majority leads by at least
  $22\sqrt{3n\log n}$ nodes (a fraction $\frac12 + O(\sqrt{\log n/n})$) and $\log n \ge 40$,
  all nodes adopt it within at most $390\log n$ rounds with probability at least
  $1 - 429\log n/n$.
\end{corollary}
\begin{proof}
  \leanok
  Theorem~\ref{thm:3_8} with $k = 2$ and $\lambda = 3$: the majority color is the unique
  plurality, holds more than $n/2 \ge n/3$ nodes, and its bias is the gap; the conclusion
  transfers to the binary process through Theorem~\ref{thm:binary}.
\end{proof}

\subsection{Two colors from any configuration}\label{sec:any_start}

This part goes beyond the paper: it follows the symmetry-breaking argument for the binary
median dynamics of Doerr, Goldberg, Minder, Sauerwald and Scheideler (SPAA 2011), as described in
Becchetti, Clementi, Natale, \emph{Consensus dynamics: an overview} (SIGACT News 2020, \S4,
Case~3), here for binary 3-majority. Let $I$ be the set of nodes with opinion $1$ and
$s = |I| - (n - |I|)$ the gap. After one round, $s' = 2\sum_v Y_v - n$ for independent
$\{0,1\}$-valued $Y_v$ with mean $p = 3x^2 - 2x^3$, $x = |I|/n$, so
$\mathbb{E}s' = s\,(3/2 - s^2/(2n^2))$.

\begin{lemma}[Jump near balance]\label{lem:jump_near_balance}
  \lean{Plurality.gap, Plurality.binKernel, Plurality.avg_pz_zero_one, Plurality.jump_near_balance}
  \leanok
  \uses{thm:binary}
  If $n \ge 5$ and $|s| \le 4\sqrt n/25$, then $|s'| \ge \sqrt n/5$ with probability at
  least $9/64$.
\end{lemma}
\begin{proof}
  \leanok
  The centered sum $S = \sum_v (Y_v - p)$ has variance $\sigma^2 = np(1-p) \ge n/5$ and fourth
  moment at most $\sigma^2 + 3\sigma^4$, so by Paley--Zygmund $4S^2 \ge \sigma^2$ with
  probability at least $9/64$. Then $|s'| = |2S + \mathbb{E}s'| \ge \sqrt{n/5} - \frac32|s|
  \ge \sqrt n/5$.
\end{proof}

\begin{lemma}[Growth away from balance]\label{lem:growth_far}
  \lean{Plurality.mean_ge, Plurality.growth_far}
  \leanok
  \uses{thm:binary}
  If $0 \le s \le n/2$ and $\lambda \ge 0$, then $s' > \frac{11}{8}s - 2\lambda$ with
  probability at least $1 - e^{-2\lambda^2/n}$.
\end{lemma}
\begin{proof}
  \leanok
  $\mathbb{E}s' \ge \frac{11}{8}s$ for $0 \le s \le n/2$, and Hoeffding's inequality bounds the
  lower tail of $\sum_v Y_v$.
\end{proof}

\begin{theorem}[Symmetry breaking]\label{thm:symmetry_breaking}
  \lean{Plurality.gapUnits, Plurality.units_near, Plurality.units_far, Plurality.units_grow, Plurality.hitProb_mono, Plurality.majority3_symmetry_breaking}
  \leanok
  \uses{lem:jump_near_balance, lem:growth_far}
  There is $C > 0$ such that, for $\log n \ge 40$, from any configuration the gap reaches
  $22\sqrt{3n\log n}$ in absolute value within any $t \ge C\log n$ rounds with probability at
  least $1 - 1/n$.
\end{theorem}
\begin{proof}
  \leanok
  Apply the hitting-time bound of Doerr et al.\ (their Claim~2.9,
  \texttt{Dynamics.Kernel.drift\_hitting\_log}) to $X = \lfloor |s|/(\sqrt n/100)\rfloor$ with
  $q = n$, $c_1 = 5/4$, $c_2 = 1/320000$, $c_3 = 9/64$, $c_4 = 1000$, $c_6 = 1$; by the
  symmetry between the opinions assume $s \ge 0$. If $X = 0$, Lemma~\ref{lem:jump_near_balance}
  gives $X' \ge 20 \ge 1$ with probability $9/64$. If $1 \le X < 16$, the same lemma gives
  $X' \ge 20 \ge \frac54 X$ with probability $9/64 \ge 1 - e^{-X/320000}$. If
  $16 \le X < 1000\log n$, then $s \le n/2$, and Lemma~\ref{lem:growth_far} with
  $\lambda = X\sqrt n/3200$ gives $s' > \frac{21}{16}X\sqrt n/100$, hence $X' \ge \frac54 X$,
  except with probability $e^{-X^2/5120000} \le e^{-X/320000}$. Reaching $X \ge 1000\log n$
  means $|s| \ge 10\sqrt n\log n \ge 22\sqrt{3n\log n}$.
\end{proof}

\begin{lemma}[Hit, then absorb]\label{lem:hit_absorb}
  \lean{Plurality.hitProb_le_one, Plurality.hitProb_nonneg, Plurality.hitProb_sub_le_event, Plurality.event_amplify}
  \leanok
  Let $P$ be an absorbing event of a finite Markov chain. If from every state of $B$ the chain is
  in $P$ after $T$ steps with probability at least $1 - \varepsilon$, then from any state it is in
  $P$ at time $t + T$ with probability at least $\Pr(T_B \le t) - \varepsilon$. If moreover from
  every state the chain is in $P$ after $T_1$ steps with probability at least $1 - \varepsilon$,
  and from $a$ after $T_0$ steps with probability at least $1 - \delta$, then from $a$ it is in
  $P$ after $T_0 + T_1$ steps with probability at least $1 - \delta\varepsilon$.
\end{lemma}
\begin{proof}
  \leanok
  The first claim by induction on $t$, through the recursion of $\Pr(T_B \le t)$ and the
  monotonicity in time of an absorbing event; the second by the Markov property at time $T_0$.
\end{proof}

\begin{theorem}[Two opinions from any configuration]\label{thm:any_start}
  \lean{Plurality.step_compl, Plurality.consensus_absorbing, Plurality.event_consensus_of_gap, Plurality.majority3_any_start, Plurality.majority3_any_start_whp}
  \leanok
  \uses{thm:symmetry_breaking, cor:binary_vanishing, lem:hit_absorb}
  There is $C > 0$ such that, for $\log n \ge 40$, from any configuration (in particular from a
  perfectly balanced one) all nodes hold the same opinion after any $T \ge C\log n$ rounds with
  probability at least $1 - C\log n/n$; and, for $\log n \ge C$, with probability at least
  $1 - 1/n$.
\end{theorem}
\begin{proof}
  \leanok
  Theorem~\ref{thm:symmetry_breaking}, then Corollary~\ref{cor:binary_vanishing} applied to the
  set of nodes of opinion $1$ or to its complement, composed by Lemma~\ref{lem:hit_absorb};
  consensus on either opinion is absorbing, so further rounds keep it. For the second bound, two
  consecutive blocks each fail with probability at most $\varepsilon = C_0\log n/n$, and
  $\varepsilon^2 \le 1/n$.
\end{proof}

\section{Lower bounds}\label{sec:lower}

\subsection{Lower bound for 3-majority}

\begin{lemma}[Escaping a moving target]\label{lem:escape}
  \lean{Dynamics.expList_escape, Dynamics.expList_sum}
  \leanok
  If a round-based process starts in $G_0$ and, for every $t < T$, one round
  from any state of $G_t$ misses $G_{t+1}$ with probability at most $p$, then
  after $T$ rounds it lies outside $G_T$ with probability at most $Tp$.
\end{lemma}

\begin{proof}
  \leanok
  Induction on $T$, conditioning on the first round.
\end{proof}

\begin{lemma}[Lemma 4.1]\label{lem:4_1}
  \lean{Plurality.lemma_4_1, Plurality.mu_le_of_count_le, Plurality.sq_div_le_sum_sq}
  \leanok
  \uses{lem:2_1, lem:tails}
  Let $k\sqrt{n\log n} \le b \le n/k$. If $c_j \le n/k + b$, then
  $\Pr(C_{j,t+1} \ge n/k + (1 + 3/k)b \mid C_t = c) \le 1/n^2$.
\end{lemma}

\begin{proof}
  \leanok
  By Cauchy--Schwarz $\sum_h c_h^2 \ge n^2/k$, so
  $\mu_j \le c_j(1 + c_j/n - 1/k)$. Writing $c_j = n/k + a$ with $a \le b$,
  this is at most $n/k + (1+1/k)b + b^2/n \le n/k + (1+2/k)b$. Hoeffding's
  inequality with deviation $b/k \ge \sqrt{n\log n}$ gives $1/n^2$.
\end{proof}

\begin{theorem}[Theorem 4.2, lower bound for 3-majority]\label{thm:4_2}
  \lean{Plurality.theorem_4_2, Plurality.theorem_4_2_log, Plurality.color_stays_small, Plurality.balancedSet}
  \leanok
  \uses{lem:4_1, lem:escape}
  Let $k \ge 3$ and suppose every color starts with at most $n/k + b$ nodes,
  where $b \ge k\sqrt{n\log n}$. If $b(1 + 3/k)^T \le n/k$, in particular if
  $T \le \frac k3 \log\frac{n}{kb}$, then the coloring after $T$ rounds is
  monochromatic with probability at most $kT/n^2$.
\end{theorem}

\begin{proof}
  \leanok
  For each color apply Lemma~\ref{lem:escape} to the thresholds
  $n/k + b(1+3/k)^t$, with Lemma~\ref{lem:4_1} for each round. A
  monochromatic coloring has a color with $n > 2n/k$ nodes, above its final
  threshold; sum over the $k$ colors.
\end{proof}

\begin{remark}
  With $b = \max\{(n/k)^{1-\varepsilon}, k\sqrt{n\log n}\}$ this gives
  $\Omega(k\log(n/(kb)))$ rounds, which is $\Omega(k\log n)$ for
  $k \le n^{1/4-\delta}$. The paper states $\Omega(k\log n)$ for all
  $k \le (n/\log n)^{1/4}$, but at that end $n/k = k\sqrt{n\log n}$ and the
  argument gives nothing. The paper's proof also chooses the thresholds after
  seeing the trajectory; the union bound needs the fixed schedule used here.
\end{remark}

\subsection{Which 3-input rules solve plurality consensus}

\begin{definition}[Rules, $\Delta$, $\delta$ and solvers]\label{def:good_rules}
  \lean{Plurality.Rule, Plurality.Conservative, Plurality.ClearMajority, Plurality.majCount, Plurality.deltaCount, Plurality.UniformRule, Plurality.Solver}
  \leanok
  \uses{def:rule, def:kernel}
  A rule is \emph{conservative} if it returns one of its inputs. It has the
  \emph{clear-majority} property if it returns the majority color whenever two
  inputs agree. For colors $r \ne b$, $\Delta_r$ counts the orderings of
  $(r,r,b)$ on which the rule returns $r$. For distinct $r,g,b$, $\delta_c$
  counts the orderings of $(r,g,b)$ on which it returns $c$; the rule is
  \emph{uniform} if $\delta_r = 2$ for all distinct triples. It is an
  \emph{$(s,\varepsilon)$-solver} if from every coloring whose unique plurality
  color $m$ has bias at least $s$, consensus on $m$ is eventually reached with
  probability at least $1-\varepsilon$.
\end{definition}

\begin{lemma}[Consensus needs survival]\label{lem:alive}
  \lean{Plurality.mono_le_alive, Plurality.count_runWith_eq_zero}
  \leanok
  \uses{def:good_rules}
  Under a conservative rule a color without nodes never returns. Hence for all
  $T, T_0$, the probability of consensus on $m$ at time $T$ is at most the
  probability that $m$ has a node at time $T_0$.
\end{lemma}

\begin{theorem}[Theorem 4.8(a), clear majority]\label{thm:4_8_a}
  \lean{Plurality.theorem_4_8_a, Plurality.not_solver_of_drift, Plurality.expected_red_le, Plurality.expected_red_le_iter}
  \leanok
  \uses{def:good_rules, lem:adopt}
  Let $8 \mid n$. With only colors $r, b$ present and $x$ the fraction of $r$,
  a node adopts $r$ with probability $p(x)$, where
  $p(x) - x = x(1-x)\bigl(x(\Delta_r - 2) + (1-x)(2 - \Delta_b)\bigr)$. If
  $\Delta_r \le 2 \le \Delta_b$, the number of $r$ nodes is a supermartingale,
  so from $5n/8$ nodes of $r$ (bias $n/4$) consensus on $r$ has probability at
  most $5/8$ at every time: the rule is not an $(n/4, 1/4)$-solver.
  Consequently, for an $(n/4, 1/4)$-solver and every pair $r \ne b$, either
  $\Delta_r = \Delta_b = 3$ or $\Delta_r, \Delta_b \le 1$.
\end{theorem}

\begin{proof}
  \leanok
  $\E[X_T] \le X_0$ by induction on $T$, and $\Pr(X_T = n) \le \E[X_T]/n$.
\end{proof}

\begin{remark}
  The paper claims that an $(n/4,1/4)$-solver has the clear-majority
  property, i.e.\ that $\Delta_r, \Delta_b \le 1$ is also impossible. Its
  proof uses the supermartingale property, but checks it only when red is the
  majority; for $\Delta_r, \Delta_b \le 1$ it fails when red is the minority,
  because the process is pulled to the interior point
  $x^* = (2 - \Delta_b)/(4 - \Delta_r - \Delta_b)$. That case is not
  formalized.
\end{remark}

\begin{lemma}[Adopting $r$ among three colors]\label{lem:adopt3}
  \lean{Plurality.adopt_r_le, Plurality.avg_adopt_r_le, Plurality.expected_count_le, Plurality.jump_prob_le}
  \leanok
  \uses{def:good_rules, lem:adopt, lem:tails}
  Let the rule be conservative with the clear-majority property, let only
  $r, g, b$ be present and $\delta_r \le 1$. Then
  $p(r) \le 3x^2 - 2x^3 + \delta_r x_r x_g x_b \le x(3x - 2x^2 + (1-x)^2/4)$,
  which is at most $0.97x$ for $x \le 2/5$; and from $x \le 2/5$ one round
  ends above $2n/5$ with probability at most $e^{-9n/31250}$.
\end{lemma}

\begin{lemma}[$r$ dies out]\label{lem:extinct}
  \lean{Plurality.frozenStep, Plurality.potential, Plurality.frozen_potential_le, Plurality.frozen_stop_le, Plurality.foldl_frozen_eq, Plurality.alive_le}
  \leanok
  \uses{lem:adopt3, lem:escape}
  Under the hypotheses of Lemma~\ref{lem:adopt3}, starting with at most $2n/5$
  nodes of $r$, the probability that $r$ still has a node after $T$ rounds is
  at most $0.97^T c_r + T e^{-9n/31250}$.
\end{lemma}

\begin{proof}
  \leanok
  Consider the chain frozen once $r$ exceeds $2n/5$. On it,
  $\Phi = C_r\,\mathbb 1[C_r \le 2n/5]$ contracts by $0.97$ in every state,
  it is frozen within $T$ rounds with probability at most $Te^{-9n/31250}$,
  and it agrees with the real chain until it is frozen.
\end{proof}

\begin{theorem}[Theorem 4.8(b), uniform property]\label{thm:4_8_b}
  \lean{Plurality.theorem_4_8_b, Plurality.not_solver_of_nonuniform, Plurality.extinctRounds, Plurality.deltaCount_sum}
  \leanok
  \uses{lem:extinct, lem:alive}
  Let $60 \mid n$ and $\log n \ge 20$. A rule with the clear-majority property
  that is an $(n/20, 1/4)$-solver has the uniform property.
\end{theorem}

\begin{proof}
  \leanok
  If a triple is not uniform, some color $r$ of it has $\delta_r \le 1$,
  since the three counts sum to $6$. Start from $23n/60$ nodes of $r$, $n/3$
  of $g$ and $17n/60$ of $b$: $r$ is the unique plurality color with bias
  $n/20$. By Lemmas~\ref{lem:extinct} and~\ref{lem:alive}, with
  $T_0 = \lceil\log(8n)/\log(100/97)\rceil$, consensus on $r$ has probability
  at most $1/8 + 1/8$ at every time.
\end{proof}

\begin{remark}
  The paper's proof (Lemma 4.10) is a sketch that treats $\delta_r = 1$ with
  $(\delta_g,\delta_b) \in \{(3,2),(4,1)\}$ and asserts that $r$ ``will not
  converge'' by iterating a one-round estimate. The proof here covers every
  non-uniform triple and makes the iteration rigorous.
\end{remark}

\subsection{A lower bound for \texorpdfstring{$h$}{h}-plurality}

\begin{definition}[$h$-plurality]\label{def:hplur}
  \lean{Plurality.HSample, Plurality.HRound, Plurality.mult, Plurality.hplur, Plurality.hstep, Plurality.hrun}
  \leanok
  \uses{def:config}
  Every node samples $h$ nodes and a uniformly random ordering $\sigma$ of the
  $h$ sample positions, and adopts the color of the position of smallest
  $\sigma$-rank among those whose color is most frequent. All tied colors
  appear equally often, so each is chosen with the same probability.
\end{definition}

\begin{lemma}[Adoption under $h$-plurality]\label{lem:hadopt}
  \lean{Plurality.hplur_first, Plurality.avg_hadopt_le, Plurality.avg_eval_two, Plurality.avg_prod_iter}
  \leanok
  \uses{def:hplur}
  A node adopts $j$ with probability at most $x_j + \frac{h^2}{2}x_j^2$.
\end{lemma}

\begin{proof}
  \leanok
  If $j$ wins and is seen exactly once, all $h$ sampled colors are distinct
  and $j$ is at the position ranked first, which has probability $x_j$.
  Otherwise $j$ is seen at least twice, and
  $\Pr(N_j \ge 2) \le \E[N_j(N_j-1)]/2 = \binom h2 x_j^2$.
\end{proof}

\begin{lemma}[Lemma 4.11]\label{lem:4_11}
  \lean{Plurality.lemma_4_11, Plurality.lemma_4_11_paper, Plurality.exp_bound_le}
  \leanok
  \uses{lem:hadopt, lem:tails}
  If $c_j \le a \le 2n/k$, then
  $\Pr(C_{j,t+1} \ge a(1 + 2h^2/k)) \le \exp(-2(ah^2/k)^2/n)$, which is at
  most $e^{-2nh^4/k^4}$ when $a \ge n/k$.
\end{lemma}

\begin{remark}
  The paper states the factor $1 + h^2/k$ but its proof gives
  $1 + 2h^2/k$, which is also what Theorem~\ref{thm:4_12} uses. Allowing
  $c_j \le a$ instead of $c_j = a$ covers colors below $n/k$, which the
  paper's proof of Theorem~\ref{thm:4_12} needs but the lemma excludes.
\end{remark}

\begin{theorem}[Theorem 4.12, lower bound for $h$-plurality]\label{thm:4_12}
  \lean{Plurality.theorem_4_12, Plurality.theorem_4_12_log, Plurality.hBalancedSet}
  \leanok
  \uses{lem:4_11, lem:escape}
  Let $k \ge 3$. If every color starts with at most $\frac32\cdot\frac nk$
  nodes and $T \le \frac{k}{2h^2}\log\frac43$, then after $T$ rounds of
  $h$-plurality the coloring is monochromatic with probability at most
  $kT e^{-2nh^4/k^4}$.
\end{theorem}

\begin{proof}
  \leanok
  As for Theorem~\ref{thm:4_2}, with thresholds
  $\frac32\cdot\frac nk(1 + 2h^2/k)^t \le 2n/k$.
\end{proof}
